For the definite count, write \(R=\mathcal O_E\) and
let \(W\) denote the definite Hermitian space of rank \(k\) obtained
from the common target row. For a full integral lattice \(A\) in a
nondegenerate subspace, write \(A^\#\) for its Hermitian dual
in that subspace and \(\Delta(A)=|A^\#/A|\). In particular,
\(\Delta(0)=1\). Fix the genus \(\mathcal G\) of
Proposition~\ref{prop:flag-row-mass}, with \(T=T_u\). Thus
\[
     L^\#/L\simeq R/TR,\qquad D=\Delta(L)=T^{2f}
          \qquad(L\in\mathcal G),\qquad f=[F:\mathbf Q].
\]

An \emph{embedded datum of index \(x\)} over \(L\) consists of
a full integral lattice \(B\subset W\) of discriminant order \(D\)
and the inverse image \(\Gamma\subset B\times L\) of the graph
of a bijective isometry \(B/gB\to L/gL\) of the forms modulo \(g\).
Its diagonal lattice must satisfy
\[
 M_\Gamma=\{z\in W:(z,z)\in\Gamma\},\qquad
 [L:M_\Gamma]=x.
\]

The quotient \(L/M_\Gamma\) must be generated by at most \(s=m/2-1\)
elements over \(R\).
Write \(\mathcal T(L;x,x')\) for the set of triples consisting of
two such data \(d=(B,\Gamma)\), \(d'=(B',\Gamma')\), of indices
\(x,x'\), and a rational unitary isometry \(\eta:W\to W\) such that
\[
          \eta B'=B,\qquad (\eta\times1)\Gamma'=\Gamma .
\]

The lattices and isometries here are counted individually inside
the one target space. Normalization by target automorphisms occurs
only when averaging over the genus.

\begin{auxlemma}[Compatible embedded data in the definite genus]
\label{lem:zero-gram-definite-mean}
Let \(k\ge3\), \(g\ge2\), and \(x,x'\ge1\) be integers, and let
\(E/F\), \(T=T_u\) and \(\mathcal G\) be as in
Proposition~\ref{prop:flag-row-mass}. Put
\[
 \begin{gathered}
   c=1+\log_2(2f),\qquad v=\mathrm{rd}(E)/2\ge1,\qquad
   X=xx',\qquad Y=XD^2,\\
   M_b=\bigl(20f(2v)^2\bigr)^{2f(6+c)},\qquad
   \mathcal A=M_b(2v)^f,\\
   U_*=(32/g)^{kf}(4kf)^{4f}Y^7,\qquad
   V_*=(2/g)^{kf}(4kf)^{4f}DX .
 \end{gathered}
\]
For the fields of \eqref{eq:host-character-compositum},
the quadratic subfield \(\mathbf Q(\sqrt{-d})\) has discriminant
of absolute value \(d\equiv7\bmod16\) by
Proposition~\ref{prop:quadratic-prime-pair}. The discriminant tower formula
gives \(2v=\mathrm{rd}(E)\ge\sqrt d\).
For the probability proportional to \(1/|\operatorname{Aut}(L)|\)
on isometry classes in \(\mathcal G\),
\begin{equation}\label{eq:zero-gram-pair-mean}
   \big\langle |\mathcal T(L;x,x')|\big\rangle_{\mathcal G}
   \le X^{k/2}Y^{6s+10+4c}
      \left(1+\mathcal A\sum_{h=1}^{k-1}U_*^h+V_*^k\right).
\end{equation}
The bound applies when the two quotients have different minimal
numbers of generators, each at most \(s\).
\end{auxlemma}

\begin{proof}
\proofstep{Data with a prescribed diagonal lattice.}
Take \(M=M_\Gamma\) and \(u_0=2s+1\). The inclusions
\[
 M\subset L\subset L^\#\subset M^\#,\qquad
 |M^\#/M|=Dx^2,\qquad |M^\#/L|=Dx
\]
give \(M^\#/M\) at most \(u_0\) generators. Indeed the two outside
quotients are dual up to conjugation and the middle quotient is cyclic.
Every submodule of this finite module has the same generator bound
by elementary divisors locally and the Chinese remainder theorem.
Since \(M\subset B\subset M^\#\), a list of \(u_0\) elements
of \(M^\#/M\) suffices to specify \(B/M\). Once \(B\) is chosen,
fix lifts \(b_i\) of a generating list for \(B/M\). If \(l_i\)
represents the image of \(b_i\) modulo \(gL\), testing the isometry
on \(b_i\) and each \(z\in M\) gives
\[
 \langle b_i-l_i,z\rangle\in gR,\qquad
             (b_i-l_i)/g\pmod L\ \in M^\#/L .
\]

These classes determine the images of the \(b_i\), and the graph
already sends each \(z\in M\) to \(z\bmod gL\). The number of
possible data above \(M\) is consequently at most
\[
        (Dx^2)^{u_0}(Dx)^{u_0}\le(Dx)^{3u_0}.
\]

\proofstep{The common fixed space and its discriminant.}
Write \(M'=M_{\Gamma'}\), \(U=\ker(1-\eta)\), \(h=k-\dim_E U\).
On \(U\) the two graph conditions coincide. Put
\[
 M_0=M\cap U=M'\cap U,\qquad a=\Delta(L\cap U),\qquad
 y=[L\cap U:M_0]\le\min(x,x').
\]

Let \(P:W\to U^\perp\) be orthogonal projection. For \(z'\in M'\)
and \(z\in M\), the pairs \((\eta z',z')\) and \((z,z)\) belong
to \(\Gamma\). Comparing their pairings in its two coordinates gives
\[
    C(PM')\subseteq(PM)^\#,\qquad
                C=(1-\eta)|_{U^\perp}/g .
\]

Definiteness makes \(C\) invertible on \(U^\perp\). To take indices
also for these projected lattices, let \(\delta(A)\) be the real
covolume in the product of the \(f\) complex Hermitian spaces,
divided by \(v^{f\operatorname{rank}_E(EA)}\). This divides out
the covolume \(2^{-f}\sqrt{|\operatorname{disc}(E)|}=v^f\)
of \(R\) in one Hermitian line. The determinant formula gives
\[
  \delta(A^\#)=\delta(A)^{-1},\qquad
  \delta(M)=\delta(M_0)\delta(PM),\qquad
  \delta(M')=\delta(M_0)\delta(PM').
\]

The last two identities follow by projecting a real fundamental
parallelepiped along \(U\). For integral \(A\),
\(\delta(A)^2=\Delta(A)\). At each complex embedding, every
eigenvalue of \(\eta\) has absolute value one. With \(q_0=2/g\)
the projected inclusion therefore yields
\begin{equation}\label{eq:zero-gram-fixed-discriminant}
  ay^2=\delta(M_0)^2
   \le |N_{E/\mathbf Q}\det C|\,\delta(M)\delta(M')
   \le q_0^{2fh}DX .
\end{equation}

In particular, only integers \(1\le a\le DX\) can occur.

\proofstep{Triple count for a specified space \(U\).}
A rank-\(n\) projective \(R\)-lattice has at most \(I^{n+c}\)
sublattices of index \(I\), by the prime-by-prime chain count in
Proposition~\ref{prop:primitive-flag-means}. Thus for a fixed \(y\)
the number of \(M_0\) is at most
\[
   y^{k-h+c}\le (q_0^{2fh}DX/a)^{k/2}X^c.
\]
There are at most \(X\) choices for \(y\). Given \(M_0\), the image
of \(M\) in \(L/(L\cap U)\) has index \(x/y\), hence at most
\((x/y)^{h+c}\) possibilities. For a fixed image \(J\), lifts form
a torsor under \(\operatorname{Hom}_R(J,(L\cap U)/M_0)\).
Locally \(J\) is free of rank \(h\), so
\[
 |\operatorname{Hom}_R(J,(L\cap U)/M_0)|=y^h,\qquad
             (x/y)^{h+c} y^h\le x^{2h+c}.
\]
This bounds the choices for \(M\). Replacing \(x\) by \(x'\) gives
the corresponding bound for \(M'\).

With both data fixed, all compatible \(\eta\) that fix \(U\)
pointwise lie in one coset of the group of automorphisms of \(B\)
fixing \(U\) pointwise. This finite group acts faithfully on
\(B/(B\cap U)\), since orthogonality identifies its rational
action with the action on \(U^\perp\). For \(h>0\), put \(n=2fh\).
By Minkowski's theorem \cite[\S1.1]{serre-finite-orders}, its
order divides \((2n)!\). Indeed the bound on the exponent at
each prime \(p\) satisfies
\[
 \sum_{i\ge0}\left\lfloor\frac{n}{p^i(p-1)}\right\rfloor
       \le \sum_{i\ge1}\left\lfloor\frac{2n}{p^i}\right\rfloor
       =\operatorname{ord}_p((2n)!) .
\]
Consequently there are at most \((4hf)^{4hf}\) isometries, with
the value \(1\) when \(h=0\). Multiplying the bounds for the
intersection, its two lifts, the two graphs and the isometry gives
the following count over a specified \(U\):
\begin{equation}\label{eq:zero-gram-one-space}
  (q_0^{2fh}DX/a)^{k/2}
        X^{2h+2c+1}Y^{3u_0}(4hf)^{4hf}.
\end{equation}

\proofstep{The mean number of fixed spaces in intermediate dimension.}
For \(1\le h<k\), write \(z=k-h\), \(H_0=aD\) and
\(\lambda=\log_2(2f)=c-1\). Such spaces \(U\) of discriminant order
\(a\) correspond bijectively to their primitive lattices \(L\cap U\).
We need the constant in the gluing calculation of
Proposition~\ref{prop:primitive-flag-means}. Write \(m_i\) for
the standard unimodular mass in positive rank \(i\). The complementary
lattice \(L\cap U^\perp\) has discriminant order \(b\le H_0\):
the gluing group injects into both discriminant modules. In fact,
\[
    H=L/((L\cap U)\oplus(L\cap U^\perp)),\qquad
    ab=D|H|^2,\qquad |H|\le a .
\]

For the rational isometry type of \(U\), signatures are fixed and
at the finite places the determinant modulo norms classifies the
Hermitian space \cite{landherr}. At split places the class is trivial. At unramified
inert places it is determined by valuation parity and is trivial
away from primes over \(a\). At an inert prime \(\mathfrak P\)
an integral basis with Gram matrix \(H_{\mathfrak P}\) gives
\[
 \operatorname{ord}_{\mathfrak P}\det H_{\mathfrak P}
 =\operatorname{length}_{R_{\mathfrak P}}
     ((L\cap U)^\#/(L\cap U))_{\mathfrak P}.
\]

Thus these lengths also specify the determinant parities. The prime lists for those modules, of
total order \(a\), number at most \(a^\lambda\) by the index
count in Proposition~\ref{prop:primitive-flag-means}.
Local-to-global classification and Witt cancellation determine the
complementary rational space for each choice. Allow \(a^{1+\lambda}\)
choices. For each type the two fixed-discriminant mass bounds and
the sum over \(1\le b\le H_0\) cost at most
\[
          m_z a^{(z+\lambda)/2},\qquad
          H_0 m_h H_0^{(h+\lambda)/2}.
\]

To glue the two lattices, choose a subgroup of the direct sum of
their discriminant modules which injects into each factor. Its
projection into the rank-\(h\) discriminant module requires at most
\(h\) generators locally, hence globally. Even allowing
\((aH_0)^{2h+1}\) choices of subgroup suffices. Orbit--stabilizer
divides this count by the two automorphism orders before summing.
With \(N_{h,a}(L)=\#\{U:\dim U=z,\ \Delta(L\cap U)=a\}\)
the resulting unnormalized mass is at most
\[
 m_zm_h\,a^{z/2}H_0^{(9h+8+4\lambda)/2}.
\]

This calculation uses positive ranks \(z,h\), so includes \(h=1\)
and \(z=1\). By Proposition~\ref{prop:flag-row-mass},
the normalizing mass is
\[
     m_k\big((T^k-1)/(T-1)\big)^f\ge m_kD^{(k-1)/2}.
\]

Since \(a^{z/2} D^{-(k-1)/2}=(a/D)^{k/2}D^{1/2}a^{-h/2}\),
division gives
\[
 \langle N_{h,a}\rangle_{\mathcal G}
   \le \frac{m_zm_h}{m_k}(a/D)^{k/2}(aD)^{5h+5+2c}.
\]

It remains to control the rank dependence of the mass ratio. Put \(r=\min(z,h)\), \(t=k-r\) and \(R_i=m_i/m_{i-1}\)
for \(i\ge2\). Shimura's adjacent-rank formula as used in
Proposition~\ref{prop:primitive-flag-means} gives
\[
 R_i=\big[(i-1)!(v/\pi)^i\big]^f(2v)^{-f/2}
        L_F(i,\chi_{E/F}^i),\qquad
 2^{-f}\le L_F(i,\chi_{E/F}^i)\le2^f .
\]
The bounds follow from the Euler product and
\(\zeta(i)^f\le\zeta(2)^f<2^f\). The rank-one estimate there,
\(m_1\le2^fh_E\), and its Minkowski ideal count
\(h_E\le(2v)^{f(1+\lambda)}\) give \(m_1\le M_b\).
Using \(v\ge1\), factorial monotonicity, and \(r+t=k\),
\[
 \begin{aligned}
  \frac{m_zm_h}{m_k}
     &=\frac{m_1}{R_{t+1}}\prod_{i=2}^r\frac{R_i}{R_{t+i}}\\
     &\le m_1(2v)^f4^{rf}\pi^{(rt+1)f}
       \le \mathcal A\,16^{hkf}.
 \end{aligned}
\]
Indeed \(r(t+1)+1\le2hk\) and \(\pi<4\).

\proofstep{Summing the three ranges.}
For \(1\le h<k\), multiply the mean just obtained by
\eqref{eq:zero-gram-one-space}. The discriminant factors cancel
in the main exponent:
\[
  (q_0^{2fh}DX/a)^{k/2}(a/D)^{k/2}
         =q_0^{fhk}X^{k/2}.
\]
Use \(X\le Y\), \(aD\le Y\) and \(3u_0=6s+3\).
For each \(a\) the result is at most
\(X^{k/2}Y^{6s+9+4c}\mathcal A U_*^h\).
There are at most \(DX\le Y\) possible orders \(a\).

When \(h=0\), \(U=W\) and \(a=D\). Direct substitution in
\eqref{eq:zero-gram-one-space} gives at most
\(X^{k/2}Y^{6s+10+4c}\).
When \(h=k\), \(U=0\) is the sole space. Counting the two
diagonal lattices before imposing compatibility costs at most
\(X^{k+c}\). With the data and isometry bounds this gives
\[
       X^{k+c}Y^{3u_0}(4kf)^{4kf}.
\]
If any such triple occurs, \eqref{eq:zero-gram-fixed-discriminant}
with \(a=y=1\) gives \(1\le q_0^{2kf}DX\), hence
\(1\le X^{k/2}(q_0^{kf}D)^k\). Multiply by the latter factor
and collect powers to obtain at most
\(X^{k/2}Y^{6s+10+4c}V_*^k\). Adding the two endpoints and
the intermediate sum proves \eqref{eq:zero-gram-pair-mean}.
\end{proof}
